% =====================================================================
% Appendix H --- The pricing door: derivations, guard and accounting
% ---------------------------------------------------------------------
% Job: the formal backing of Sections 6.4-6.5 (subsec:pricing-capacity,
%      subsec:pricing-markup): the effective-technology and markup hooks,
%      the first-order conditions at marginal cost, the uniqueness guard,
%      the Cobb-Douglas price identity, and the external-account and GDP
%      accounting with the profit booked.
% Code: OpenCGE@src/closures/pricing/pricing.jl, OpenCGE@src/core/equilibrium.jl
%      (_mobile_market_demand, _cost_minimizing_labor), OpenCGE@src/core/
%      diagnostics.jl (gdp_components, retained_profit); the gate lives in
%      experiments/run.jl (evaluate_gates). Records: ADR-0029, ADR-0030, ADR-0031.
% =====================================================================
\section{The pricing door: derivations, guard and accounting}\label{app:pricing-door}

\subsection{Two hooks on one cost line}\label{app:pricing-hooks}
Sector $i$ produces with the CES unit cost
\begin{equation}\label{eq:app-unit-cost}
  c_i \;=\; \Bigl[\,(A^{\mathrm{eff}}_i)^{\epsilon-1}\bigl(\alpha_i\,w_i^{1-\epsilon}
     + (1-\alpha_i)\,P_i^{1-\epsilon}\bigr)\Bigr]^{\frac{1}{1-\epsilon}},
\end{equation}
where $P_i$ is the intermediate price index (CES in $\theta$ over the
purchaser prices $p_u$) and $\alpha_i$ the factor share. The two doors act on
this one line. The capacity door replaces $A_i$ by
$A^{\mathrm{eff}}_i = A_i (y_i/\lambda_i)^{-\delta_i \alpha_i}$ ($\lambda_i$ the sector's baseline gross output) and keeps zero profit,
$p_i = c_i$. The factor share in the exponent converts the value-added
elasticity $\delta_i$ into a gross-output one: at fixed input prices a
Hicks-neutral change $d\ln A^{\mathrm{eff}}_i$ raises value added per worker by
$d\ln A^{\mathrm{eff}}_i/\alpha_i$, so value-added productivity responds to own
output with elasticity $-\delta_i$. The markup door keeps $A^{\mathrm{eff}}_i = A_i$ and sets
$p_i = \mu_i c_i$ with $\mu_i = 1 + \kappa(y_i/\lambda_i - 1)$. Both are
written in $y_i/\lambda_i$, so the reference point ($y = \lambda$) is a
fixed point of either hook ($A^{\mathrm{eff}} = A$, $\mu = 1$) and the
scale-determinacy guard of the kernel, evaluated at the reference, is
unchanged. Their product
$p_i = \bigl(1 + \kappa(y_i/\lambda_i - 1)\bigr)\,c_i\bigl(A_i (y_i/\lambda_i)^{-\delta_i\alpha_i}\bigr)$
is one family with two parameters; $\delta = \kappa = 0$ reproduces the
kernel exactly (bit for bit, as the nesting cells
\texttt{pricing\_capacity-v2-GAMMA-*-d0},
\texttt{pricing\_capacity\_rows-v1-*-d0} and
\texttt{pricing\_markup-v2-*-k0} confirm against their
\texttt{matrix\_5x3-v10} twins on every labor closure).

\subsection{First-order conditions at marginal cost}\label{app:pricing-foc}
A firm with a markup buys at the purchaser prices but equates marginal revenue
product to factor cost at its \emph{marginal cost} $c_i = p_i/\mu_i$. Cost
minimization (Shephard's lemma on \eqref{eq:app-unit-cost}) gives labor demand
\begin{equation}\label{eq:app-labour}
  L_i \;=\; c_i^{\epsilon}\,(A^{\mathrm{eff}}_i)^{\epsilon-1}\,\alpha_i\,y_i\,w_i^{-\epsilon},
\end{equation}
and the demand for input $u$ of user $i$ is
\begin{equation}\label{eq:app-intermediate}
  x_{ui} \;=\; p_u^{-\theta}\,\Omega_{iu}\;
     c_i^{\epsilon}\,(A^{\mathrm{eff}}_i)^{\epsilon-1}\,P_i^{\theta-\epsilon}\,a_i\,y_i ,
\end{equation}
with $a_i$ the domestic intermediate bill per unit of output. The outer
$p_u^{-\theta}$ is the price of the demanded input and stays the purchaser
price; the inner $c_i^{\epsilon}$ is the user's marginal cost. The same factor
$k_i = c_i^{\epsilon}(A^{\mathrm{eff}}_i)^{\epsilon-1}P_i^{1-\epsilon}$ values
the two leaks of the purchaser-price intermediate bill, imported
intermediates $M^{\mathrm{int}}$ and product taxes on intermediate use
$T^{\mathrm{int}}$. These are the four sites where the price enters a
first-order condition (labor demand, the inner factor of intermediate demand,
and both leak valuations); writing $p$ instead of $c$ at any of them inflates
the corresponding demand by $\mu^{\epsilon}$ while leaving the price response
untouched. The two leak sites are invisible on any synthetic fixture
($M^{\mathrm{int}} = T^{\mathrm{int}} = 0$), which is why the paper-side tests
assert the external account on the German table. Under the capacity door
$\mu \equiv 1$, so $c = p$ and only $A^{\mathrm{eff}}$ moves: the capacity door
needs no first-order-condition change.

\subsection{The Cobb--Douglas price identity}\label{app:pricing-cd}
At $\theta = \epsilon = 1$ and a pinned wage $w = 1$ the price block is
log-linear: $\log p_i = \log \mu_i - \log A^{\mathrm{eff}}_i
+ (1-\alpha_i)\sum_u \Omega_{iu}\log p_u$, hence
\begin{equation}\label{eq:app-cd}
  \log p \;=\; (I - B)^{-1}\bigl(\log \mu - \log A^{\mathrm{eff}}\bigr),
  \qquad B = \operatorname{diag}(1-\alpha)\,\Omega .
\end{equation}
A uniform constant markup $\bar\mu$ therefore raises prices by
$\bar\mu^{(I-B)^{-1}\mathbf 1}$, not by $\bar\mu$: the markup cascades through
the intermediate network, by a factor of about $1.8$ on the German table. The
kernel tests assert \eqref{eq:app-cd} for both doors.

\subsection{The uniqueness guard}\label{app:pricing-guard}
With $\delta_i < 0$ a sector's unit cost falls as its output rises, a positive
feedback that can admit several equilibria. Every priced cell is reached by a
deterministic continuation in its door parameter (four equal steps from the
door-free solution) and must pass a uniqueness gate: re-solved from four
starting points (the default initial point, the door-free solution, and the
solution with all quantities scaled by $0.9$ and $1.1$), at least two must
converge and every converged start must lie within $10^{-6}$ of the solution,
and the Jacobian of the full residual must have full rank. The measured
spreads are at most $6 \cdot 10^{-9}$ at condition numbers up to $10^{8}$
(the no-reallocation closure, whose sectoral wages enlarge the system); a
distinct equilibrium would sit about $10^{-2}$ away. All 180 priced cells pass,
on every labor closure.

\subsection{Accounting with the profit booked}\label{app:pricing-accounting}
The profit bill is $\Pi$ of equation~\eqref{eq:profit-bill}. With retention
rate $s_\pi$ the household spends
$E = (1-\tau)\,w\textstyle\sum_i L_i + F + (1-s_\pi)\Pi$, and the
external-account identity of the kernel becomes
\begin{equation}\label{eq:app-account}
  sE + s_\pi \Pi + T^{\mathrm{int}} + M - (I + X) \;=\; F + B_{\mathrm{gov}},
\end{equation}
which every priced run satisfies to $10^{-11}$; dropping the $s_\pi\Pi$ term
leaves it short by exactly the retained profit. On the income side GDP is factor
income $w\sum_i L_i + \Pi$, so the companion identity between the GDP wedge and
the account gap holds unchanged. Retention ($s_\pi = 1$) needs injections in
the wage-pinned gauge: summing over sectors,
$(\mu - 1)\,c\,y + (s + m - s m)\,w L = p\,(1-m)(g^G + \mathit{exo}) + pX$, so an
economy without government demand, investment or exports has no positive
retained-profit equilibrium. The application carries all three.
%TODO (relabel): Pi collides with the consumption price index Pi(p) of Sections
% 4-5 and Appendices A-B; the price index is to be relabelled before submission
% (see the note in Section 6.2 at eq. profit-bill).

\subsection{The pricing-door profiles on every labor closure}\label{app:profiles}

The main text carries the pricing door in the profile in which the utilization
elasticity is confined to the goods sectors, because the entire direct incidence
of the program falls there and because \textcite{Kaldor1966} reads Verdoorn's
law as a law of manufacturing. Table~\ref{tab:pricing-profiles} prints the rows
this leaves out, on every labor closure: the counterfactual in which $\delta$
extends to every sector, and the $\kappa = 0.25$ markup with its profit
distributed.

The all-sector profile is worth stating explicitly, because it is what makes the
sign of the price response look like a property of the financing closure. With
$\delta$ on 36 service sectors, which receive none of the program directly,
Kaldor--Verdoorn lowers the GDP deflator except under external finance
(\texttt{pricing\_capacity-v2-GAMMA-\{F1,F2,F3\}-dm05}: $-0.20$, $-0.16$,
$+0.12$ percent), against $-0.06$ to $-0.04$ percent in all three columns when
the elasticity is confined to goods. It is a statement about where returns are assumed to live,
and it carries over to the Leontief closure ($-0.16$, $-0.13$, $+0.14$ percent;
\texttt{pricing\_capacity\_rows-v1-DELTA-\{F1,F2,F3\}-dm05}), while the
standard CGE closure lowers the deflator by 0.16 percent in every column.

The three CPI-gauge closures (no-reallocation, standard CGE, elastic supply)
keep F2 and F3 indistinguishable in every row of both panels, so the financing
the transfer neutrality behind \eqref{eq:finstrument} survives the pricing door; in panel B the
distributed profit leaves them unchanged as well, because the net external
transfer absorbs it (Section~\ref{subsec:pricing-markup}).

% =====================================================================
% Table E --- pricing-door profiles not carried in the main text (GENERATED)
% ---------------------------------------------------------------------
% Do NOT edit by hand. Regenerate with:  make_pricing_door.py
% OpenCGE tag v0.5.0 (ADR-0044); source run ids are listed below.
% Sources: pricing_capacity_rows-v2-<row>-<fin>-<dp05|dm05>,
% pricing_capacity-v3-GAMMA-<fin>-<dp05|dm05>,
% pricing_mobile-v3-ALPHA-<fin>-<dp05|dm05>.
% One panel: the all-sector capacity counterfactual; the kappa = 0.25
% distributed-profit rows were dropped from the paper on 2026-10-24
% (they repeated the invariance Section 6.2 already states).
% =====================================================================
\begin{table}[htbp]
\centering
\scriptsize
\setlength{\tabcolsep}{1.6pt}
\renewcommand{\arraystretch}{1.15}
\resizebox{\textwidth}{!}{%
\begin{tabular}{@{}l@{\hspace{6pt}}rrrrr@{\hspace{6pt}}rrrrr@{\hspace{6pt}}rrrrr@{}}
\toprule
 & \multicolumn{5}{c}{\textbf{F1}} & \multicolumn{5}{c}{\textbf{F2}} & \multicolumn{5}{c}{\textbf{F3}} \\
\cmidrule(lr){2-6}\cmidrule(lr){7-11}\cmidrule(lr){12-16}
\multicolumn{1}{c}{Labor closure and door} & \multicolumn{1}{c}{$\Delta P$} & \multicolumn{1}{c}{\shortstack{$\max$\\$|\Delta p|$}} & \multicolumn{1}{c}{$\Delta L$} & \multicolumn{1}{c}{$\Delta W_C$} & \multicolumn{1}{c}{$\Delta\Pi/Y$} & \multicolumn{1}{c}{$\Delta P$} & \multicolumn{1}{c}{\shortstack{$\max$\\$|\Delta p|$}} & \multicolumn{1}{c}{$\Delta L$} & \multicolumn{1}{c}{$\Delta W_C$} & \multicolumn{1}{c}{$\Delta\Pi/Y$} & \multicolumn{1}{c}{$\Delta P$} & \multicolumn{1}{c}{\shortstack{$\max$\\$|\Delta p|$}} & \multicolumn{1}{c}{$\Delta L$} & \multicolumn{1}{c}{$\Delta W_C$} & \multicolumn{1}{c}{$\Delta\Pi/Y$} \\
\midrule
\multicolumn{16}{@{}l}{\textit{Capacity effects on all sectors: No-reallocation}} \\
\quad $\delta = +0.5$ & $+0.74$ & $26.82$ & $0.00$ & $+0.10$ & $0.00$ & $+0.86$ & $35.47$ & $0.00$ & $-1.99$ & $0.00$ & $+0.86$ & $35.47$ & $0.00$ & $-1.99$ & $0.00$ \\
\quad $\delta = -0.5$ & $+0.48$ & $16.39$ & $0.00$ & $+0.06$ & $0.00$ & $+0.53$ & $19.52$ & $0.00$ & $-1.99$ & $0.00$ & $+0.53$ & $19.52$ & $0.00$ & $-1.99$ & $0.00$ \\
\midrule
\multicolumn{16}{@{}l}{\textit{Capacity effects on all sectors: Standard CGE}} \\
\quad $\delta = +0.5$ & $+0.12$ & $4.06$ & $0.00$ & $+0.15$ & $0.00$ & $+0.13$ & $4.27$ & $0.00$ & $-1.83$ & $0.00$ & $+0.13$ & $4.27$ & $0.00$ & $-1.83$ & $0.00$ \\
\quad $\delta = -0.5$ & $-0.16$ & $4.56$ & $0.00$ & $+0.07$ & $0.00$ & $-0.16$ & $4.47$ & $0.00$ & $-1.87$ & $0.00$ & $-0.16$ & $4.47$ & $0.00$ & $-1.87$ & $0.00$ \\
\midrule
\multicolumn{16}{@{}l}{\textit{Capacity effects on all sectors: Elastic labor supply ($\eta_s = 0.5$)}} \\
\quad $\delta = +0.5$ & $+0.38$ & $14.11$ & $+0.15$ & $+0.38$ & $0.00$ & $+0.42$ & $16.21$ & $+0.16$ & $-1.64$ & $0.00$ & $+0.42$ & $16.21$ & $+0.16$ & $-1.64$ & $0.00$ \\
\quad $\delta = -0.5$ & $+0.13$ & $4.31$ & $+0.07$ & $+0.40$ & $0.00$ & $+0.13$ & $4.54$ & $+0.08$ & $-1.57$ & $0.00$ & $+0.13$ & $4.54$ & $+0.08$ & $-1.57$ & $0.00$ \\
\midrule
\multicolumn{16}{@{}l}{\textit{Capacity effects on all sectors: Employment adjustment}} \\
\quad $\delta = +0.5$ & $+0.12$ & $4.06$ & $+0.04$ & $+0.21$ & $0.00$ & $+0.10$ & $4.22$ & $+0.35$ & $-1.30$ & $0.00$ & $-0.10$ & $3.90$ & $+2.61$ & $+2.14$ & $0.00$ \\
\quad $\delta = -0.5$ & $-0.20$ & $4.60$ & $-0.15$ & $-0.53$ & $0.00$ & $-0.16$ & $4.47$ & $-0.01$ & $-1.89$ & $0.00$ & $+0.12$ & $4.03$ & $+1.05$ & $+2.43$ & $0.00$ \\
\midrule
\multicolumn{16}{@{}l}{\textit{Capacity effects on all sectors: Leontief closure}} \\
\quad $\delta = +0.5$ & $+0.14$ & $4.52$ & $+0.02$ & $+0.21$ & $0.00$ & $+0.10$ & $4.54$ & $+0.31$ & $-1.34$ & $0.00$ & $-0.15$ & $4.21$ & $+2.45$ & $+1.89$ & $0.00$ \\
\quad $\delta = -0.5$ & $-0.16$ & $4.12$ & $-0.12$ & $-0.44$ & $0.00$ & $-0.13$ & $4.15$ & $+0.03$ & $-1.79$ & $0.00$ & $+0.14$ & $3.76$ & $+1.22$ & $+2.61$ & $0.00$ \\
\bottomrule
\end{tabular}}
\caption{The capacity door with the utilization elasticity extended to every sector ---
the all-sector counterfactual of Table~\ref{tab:pricing-kaldor} --- on all five
labor closures. Conventions as in Table~\ref{tab:pricing-markup}, except that the
welfare column is household consumption ($\Delta W_C$ of Table~\ref{tab:results-matrix}).}
\label{tab:pricing-profiles}
\end{table}

